\section{问题三：运输与中继无人机联合调度}

\subsection{连续通信判定模型}

问题三继承问题二的载荷、能量、时限及资源周转约束，并要求运输无人机在爬升、巡航、下降和物资交接全过程保持通信。本文将每个运输架次展开为分段连续三维轨迹；飞行段内经纬度和海拔按时间线性插值，交接段位置固定于服务区作业高度。

对任意通信端点 $i,j$，沿两点视线对DEM进行加密采样。若任一内部采样点地面高程不低于视线海拔，则令遮挡变量 $b_{ij}(t)=1$。传播损耗为
\[
L_{ij}(t)=32.45+20\log_{10}f+20\log_{10}D_{ij}(t)
            +L_{\rm obs}b_{ij}(t),
\]
其中 $f$ 以MHz计，$D$ 以km计。双向链路允许损耗取两个方向门限的较小值。由附件参数计算，运输机--G01直连、运输机--中继接入及中继--G01回传门限分别为122、116和126~dB。

若运输机与G01链路可用，则优先记为直连；否则仅当运输机--中继接入链路和中继--G01回传链路在同一时刻均可用时，记为中继状态。其余情况为通信中断，并作为不可行条件。

\subsection{候选悬停点与联合调度}

候选点必须位于DEM范围内、离地高度不超过300~m并满足回传链路可用。以既有三个中继点为起点，先在新运输航线上检验空间覆盖，再搜索邻域位置与高度。空间覆盖与时间可行性分开检查：即使三个点共同覆盖全部采样点，两架中继机也未必能在交付时限前完成返场、周转和再次建链。

本轮从问题二22架次、5830.82~s的方案出发。直接沿用原路线时，S009--S012和S007--S003航段出现旧驻守点未覆盖的短区间；仅调整驻守点得到的部分空间可行候选，又在资源与时限模型中不可行。最终只改变两个医疗箱的批次归属：将S002-MED-01由原Q2-003并入Q2-011，将S012-MED-01由Q2-004并入Q2-012，运输架次数仍为22。前者避免为了S002医疗箱提前执行S004任务，后者去掉S009--S012跨区航段；其余覆盖缺口通过西部中继位置调整消除。

对上述批次枚举物理可行机型和访问顺序，预计算中继需求相对区间，并在区间两端增加3~s时间缓冲。CP-SAT同时决定机型、顺序、开始时刻及三个中继窗口；运输无人机与电池分别设置容量约束，中继无人机占用区间延长至返场后300~s，从而允许两架中继机自由分配，而非固定数组中的任务归属。第一阶段最小化联合返场时刻，第二阶段保持该时间上界，最小化运输与中继总能耗；随后局部调整中继位置并再次排程。

最终为A型8架次、B型8架次、C型6架次。R02驻守西部点，R01先执行东南任务，再返场换组件执行北部任务。本方案的开始时刻、机型和部分访问顺序不同于Q2，不能把5830.82~s直接当作Q3结果。固定修正批次、位置和需求分配后的毫秒模型得到6926.152~s的最优值，但该结论不覆盖其他组批、驻守点或中继需求分配，不能宣称全局最优。

中继架次往返飞行时间沿用三阶段口径。水平能耗按巡航功率乘巡航时间计算，爬升附加能耗取计划起飞总质量；建链阶段计悬停功率，服务阶段计悬停功率与通信附加功率之和。各架次满足20\%返航余量，R02两架次之间另计300~s周转时间。

\subsection{联合调度结果}

\begin{table}[htbp]
\centering
\caption{问题三中继架次方案}
\small
\begin{tabular}{cllrrrrr}
\toprule
架次 & 中继机 & 组件 & 经度 & 纬度 & 海拔/m & 服务时段/s & 能耗/kWh\\
\midrule
R01 & R02 & EC01 & 109.2066528 & 23.0432443 & 744.723 & 2295.45--6414.61 & 1.5177\\
R02 & R01 & EC02 & 109.2663164 & 23.0128425 & 629.728 & 754.45--3331.71 & 0.9968\\
R03 & R01 & EC03 & 109.2312200 & 23.0612224 & 322.250 & 5894.82--6226.21 & 0.3770\\
\bottomrule
\end{tabular}
\end{table}

联合方案包含22个运输架次和3个中继架次，使用2架中继无人机及3组能源组件。80个货箱均按时交付，硬时限违约和普通物资迟到均为0。运输能耗为65.6721~kWh，中继能耗为2.8915~kWh，总能耗为68.5635~kWh。运输最晚返航与最后中继返场均约为6926.15~s，联合完成时间为6926.1512~s。

相对原24运输架次、6997.28~s、73.5972~kWh方案，新方案减少2个运输架次，联合完成时间缩短71.13~s（1.02\%），总能耗减少5.0337~kWh（6.84\%）。旧22架次、7970.9~s、70.4145~kWh方案也保留为历史对照，不再称为当前低能耗推荐。以上均为在统一口径下得到的可行结果，未证明3个中继架次或22个运输架次的全局最少性。

\subsection{独立检验与离散粒度}

独立程序从Q3输出重新构造80箱运输、无人机与电池时间线、中继往返飞行和能源组件时间线，并复算每个链路的地形遮挡、自由空间损耗和双向门限。0.5~s粒度复核66660个通信点，未发现中断。全部中继悬停点位于DEM范围内、离地高度不超过300~m，两架中继无人机无任务重叠，能源组件均满足返航余量和充满后复用要求。

进一步采用0.1~s时间粒度复核331961个通信点，未发现中断。该结论是离散采样验证，不能替代连续时间链路可用性的严格数学证明；时间缓冲也不等于链路功率的鲁棒裕度。下一章Q4已冻结本轮运输与中继关系，重新核算任务分区和独立资源需求。
